\newpage
\begin{center}
  \textbf{\large LỜI NÓI ĐẦU}
\end{center}

Luận văn nghiên cứu bài toán phát hiện và cảnh báo hành vi không được phép trong phân xưởng nhà máy dưới dạng phát hiện hành động không--thời gian (STAD) theo từng người, tức là tại mỗi khung hình khóa hệ thống phải định vị được từng người và gán cho người đó một hoặc nhiều nhãn hành vi kèm điểm tin cậy. Bốn hành vi mục tiêu nằm trên hai thang thời gian trái ngược nhau. Ngủ trên bàn hay gác chân lên bàn là những trạng thái kéo dài và ít chuyển động còn sử dụng điện thoại hay hành vi nguy hiểm lại diễn ra ngắn và gắn với vật thể nhỏ nên khó có một toán tử thời gian duy nhất làm tốt được cả hai.

Từ đó, luận văn đề xuất \emph{Hybrid FGTN-Mamba}, một khối temporal neck đặt sau backbone Video Swin-Tiny và trước 3D RoIAlign. Khối gồm một nhánh mô hình không gian trạng thái chọn lọc quét hai chiều để nắm phụ thuộc dài hạn với chi phí tuyến tính theo $T$ và một nhánh self-attention trong cửa sổ $W=5$ để giữ biến đổi ngắn hạn, hai nhánh được hợp nhất bằng cổng sigmoid học theo từng kênh và từng bước thời gian. Khối được áp dụng độc lập trên hai tầng đặc trưng, sau đó đặc trưng vùng người của hai tầng được nối lại (multi-scale RoI) rồi đi qua ACRN và head sigmoid đa nhãn được huấn luyện bằng focal loss kết hợp trọng số lớp.

Luận văn chỉ báo cáo các lượt chạy có nhật ký đi kèm, mỗi cấu hình một seed. Trên bộ dữ liệu phân xưởng được chia lại theo video nguồn để loại bỏ rò rỉ giữa tập huấn luyện và tập kiểm định, Hybrid FGTN-Mamba tầng 4 đạt 86,66\% \mAP{} trên ba lớp có mẫu kiểm định, cao hơn mô hình không neck 2,94 điểm (3,01 điểm tại cùng epoch), chủ yếu nhờ lớp sử dụng điện thoại. Trên tập con AVA~v2.2, khối này cũng cao hơn mô hình không neck 2,43 điểm (16,42\% so với 13,99\%) và cải thiện ở cả ba nhóm hành vi. Pipeline đầy đủ với đa tỉ lệ và ACRN, huấn luyện 3 epoch, đạt 81,78\% và chưa vượt khối đơn tầng. Phép thăm dò che người cho thấy lớp hành vi nguy hiểm được nhận dạng qua bối cảnh nên AP 100\% của lớp này không phản ánh năng lực nhận dạng hành vi. Khi chạy đầu--cuối với bộ phát hiện YOLO26-x, mô hình đạt F1 gộp 51,4\% và logic cảnh báo theo track phát hiện được 16 trong 23 tình huống kiểm định; bộ phát hiện người (recall 63\%) và lớp gác chân lên bàn là các điểm cần cải thiện trước khi triển khai.

\textbf{Từ khóa:} phát hiện hành động không--thời gian; giám sát an toàn phân xưởng; Hybrid FGTN-Mamba; mô hình không gian trạng thái chọn lọc; ACRN; Video Swin Transformer.

\clearpage
\newpage
\begin{center}
  \textbf{\large ABSTRACT}
\end{center}

This thesis addresses the detection of prohibited behaviours on a factory floor as person-level spatio-temporal action detection (STAD): at every keyframe, the system localises each person and assigns one or several behaviour labels with a confidence score. The four target behaviours span two opposite temporal scales, long and almost static states (sleeping on the table, feet on the table) and short behaviours tied to small objects (using a mobile phone, dangerous actions), so a single temporal operator struggles to serve both.

We propose \emph{Hybrid FGTN-Mamba}, a temporal neck placed after a Video Swin-Tiny backbone and before 3D RoIAlign. It combines a bidirectional selective state-space branch (long-term, cost linear in $T$) with a windowed self-attention branch ($W=5$, short-term), fused by a sigmoid gate learned per channel and per time step. The block is applied independently to two backbone stages; the person RoI features of both stages are concatenated (multi-scale RoI) and passed through an Actor-Centric Relation Network (ACRN) and a multi-label sigmoid head trained with a class-weighted focal loss.

Only runs with training and evaluation logs are reported, each with a single seed. On the private factory dataset, re-split by source video to remove train--validation leakage, the stage-4 Hybrid FGTN-Mamba reaches 86.66\% mAP over the three classes with validation samples, 2.94 points above the no-neck baseline (3.01 points at the same epoch), mostly on phone use. On a subset of AVA~v2.2 the block also improves on the no-neck baseline by 2.43 points (16.42\% vs. 13.99\%) across all three action groups. The full pipeline with multi-scale features and ACRN, trained for 3 epochs, reaches 81.78\% and does not yet beat the single-stage block. A masking probe shows that the dangerous-action class is recognised from scene context, so its 100\% AP does not measure behaviour recognition. Run end to end with a YOLO26-x person detector, the model reaches a micro F1 of 51.4\% and the track-level alert logic fires on 16 of 23 validation episodes; person detection (63\% recall) and the feet-on-table class are the main gaps before deployment.

\textbf{Keywords:} spatio-temporal action detection; factory safety monitoring; Hybrid FGTN-Mamba; selective state-space model; ACRN; Video Swin Transformer.
\clearpage
